\documentclass[12pt, a4paper]{article}

% --- Packages ---
\usepackage[margin=1in]{geometry}
\usepackage{fancyhdr}
\usepackage{hyperref}
\usepackage{xcolor}
\usepackage{listings}
\usepackage{graphicx}
\usepackage{titlesec}
\usepackage{tcolorbox}

% --- Colors & Styling ---
\definecolor{primarycolor}{RGB}{0, 83, 156}
\definecolor{codebg}{RGB}{245, 245, 245}
\definecolor{keywordcolor}{RGB}{0, 0, 255}
\definecolor{stringcolor}{RGB}{163, 21, 21}

\hypersetup{
    colorlinks=true,
    linkcolor=primarycolor,
    urlcolor=primarycolor,
    pdftitle={PyCharm \& Jupyter Setup Guide}
}

\lstset{
    backgroundcolor=\color{codebg},
    basicstyle=\ttfamily\small,
    breaklines=true,
    frame=single,
    rulecolor=\color{lightgray},
    keywordstyle=\color{keywordcolor}\bfseries,
    stringstyle=\color{stringcolor},
    showstringspaces=false,
    xleftmargin=1em,
    xrightmargin=1em,
    aboveskip=1em,
    belowskip=1em
}

% --- Header & Footer (Agreed Format) ---
\pagestyle{fancy}
\fancyhf{}
\fancyhead[R]{\textbf{ Sujeet Banerjee}}
\fancyhead[L]{\textit{Data Science \& Gen AI}}
\fancyfoot[C]{\thepage}
\renewcommand{\headrulewidth}{0.4pt}
\renewcommand{\footrulewidth}{0.4pt}

% --- Title Format ---
\titleformat{\section}{\Large\bfseries\color{primarycolor}}{\thesection}{1em}{}[\titlerule]
\titleformat{\subsection}{\large\bfseries}{\thesubsection}{1em}{}

\begin{document}

\begin{center}
    {\Huge \textbf{\color{primarycolor} PyCharm \& Jupyter Notebook}}\\[0.5em]
    {\Large \textit{Installation and Virtual Environment Setup Guide}}\\[1.5em]
\end{center}

\section{Installing PyCharm for Data Science}

PyCharm is a powerful IDE by JetBrains. For native, fully-featured Jupyter Notebook support, \textbf{PyCharm Professional} is highly recommended, though community plugins exist for the Community Edition.

\begin{enumerate}
    \item Navigate to the official JetBrains website: \url{https://www.jetbrains.com/pycharm/download/}
    \item Download the installer for your operating system (Windows, macOS, or Linux).
    \item Run the installer and follow the on-screen instructions.
    \item Launch PyCharm and select \textbf{New Project}. Choose a directory for your workspace (e.g., \texttt{capstone-lab}).
\end{enumerate}

\section{Creating a Virtual Environment (venv)}

A virtual environment isolates your project's dependencies from your global Python installation.

\begin{enumerate}
    \item install Python 3.x from (python.org), if not already present. The installation instructions are available at: [python.org]
    \item Open the \textbf{Terminal} inside PyCharm (usually located at the bottom of the window, or accessible via \texttt{Alt + F12} / \texttt{Option + F12}).
    \item Ensure you are in your project's root directory.
    \item Run the following command to create a virtual environment named \texttt{venv}:
\end{enumerate}

\begin{lstlisting}[language=bash, title={Terminal - Create venv}]
python -m venv venv
\end{lstlisting}

\begin{enumerate}
    \setcounter{enumi}{3}
    \item Activate the virtual environment. The command depends on your operating system:
\end{enumerate}

\textbf{For Windows (Command Prompt / PowerShell):}
\begin{lstlisting}[language=bash]
.\venv\Scripts\activate
\end{lstlisting}

\textbf{For macOS and Linux:}
\begin{lstlisting}[language=bash]
source venv/bin/activate
\end{lstlisting}

\textit{Note: You will know the environment is active when you see \texttt{(venv)} prepended to your terminal prompt.}

\section{Installing Dependencies from requirements.txt}

With your virtual environment active, you must install the required libraries (including Jupyter itself) using your \texttt{requirements.txt} file.

\begin{enumerate}
    \item Ensure your \texttt{requirements.txt} file is placed in the root of your project directory. 
    \item Upgrade pip to the latest version (best practice):
\begin{lstlisting}[language=bash]
python -m pip install --upgrade pip
\end{lstlisting}
    \item Install the requirements:
\begin{lstlisting}[language=bash]
pip install -r requirements.txt
\end{lstlisting}
\end{enumerate}

\begin{tcolorbox}[colback=yellow!10!white, colframe=orange!80!black, title=Crucial Requirement]
To run notebooks, your \texttt{requirements.txt} \textbf{must} include the \texttt{jupyter} or \texttt{notebook} package. 
\end{tcolorbox}

\section{Configuring PyCharm to Execute .ipynb Files}

Now that the environment is set up and dependencies are installed, you need to tell PyCharm to use this specific \texttt{venv} to run your notebooks.

\subsection{Step 1: Set the Python Interpreter}
\begin{enumerate}
    \item Go to \textbf{File $\rightarrow$ Settings} (Windows) or \textbf{PyCharm $\rightarrow$ Settings} (macOS).
    \item Navigate to \textbf{Project: [Your Project Name] $\rightarrow$ Python Interpreter}.
    \item Click the \textbf{Add Interpreter} link or the gear icon, and select \textbf{Add Local Interpreter}.
    \item Choose \textbf{Virtualenv Environment}, select \textbf{Existing environment}, and browse to the \texttt{python.exe} (Windows) or \texttt{python} (macOS/Linux) executable inside your newly created \texttt{venv} directory.
    \item Click \textbf{Apply} and \textbf{OK}.
\end{enumerate}

\subsection{Step 2: Executing the Notebook}
\begin{enumerate}
    \item In the Project Explorer, right-click your project folder and select \textbf{New $\rightarrow$ Jupyter Notebook}. Name it (e.g., \texttt{analysis.ipynb}).
    \item Open the \texttt{.ipynb} file. PyCharm will display its native notebook UI.
    \item In the top right corner of the notebook editor, look for the \textbf{Jupyter Server} dropdown.
    \item Select \textbf{Managed Server} (PyCharm will automatically start a local Jupyter server using your configured \texttt{venv} interpreter).
    \item Type a test command in the first cell, such as:
\begin{lstlisting}[language=Python]
import pandas as pd
print("Jupyter Notebook is running successfully!")
\end{lstlisting}
    \item Click the green \textbf{Play} button next to the cell or press \texttt{Shift + Enter} to execute the code. The output will appear directly beneath the cell.
\end{enumerate}

\end{document}